\chapter{Arrays}\label{chap:arrays}

\section{Design rationale}

We implemented dynamic arrays that can grow or shrink automatically. To maintain compatibility with standard C arrays and syntax, we chose to store metadata (such as the current length and the allocated size) in a hidden header located just before the array pointer. This allows the user to use the array pointer directly as a standard C array.

The implementation uses macros to provide a convenient interface for creating, resizing, and accessing array elements. We use `void*` pointers and `memcpy`/`memmove` for generic operations, but the macros cast them to the appropriate types for type safety and ease of use.

To minimize the number of memory allocations, we use an amortized allocation strategy. When the array needs to grow, we allocate more memory than strictly necessary (geometric growth). This reduces the frequency of `realloc` calls, improving performance for repeated insertions.

\section{Array interface}\label{arrays:interface}

The interface provides macros for creating arrays (\texttt{array}), getting their length (\texttt{array\_length}), pushing elements (\texttt{array\_push}), popping elements (\texttt{array\_pop}), and deleting the array (\texttt{array\_delete}).

\begin{code}
  \lstinputlisting[style=C,linerange={6-22}]{libellul/include/libellul/type/array.h}
  \caption{Interface for (possibly variable-length) arrays.}
\end{code}
